% problem-000405 4 晶体结构与晶胞
\section*{4. 晶体结构与晶胞}
\textit{下列为分问。}

\subsection*{4-1}
⾃然界中最常⻅的冰是冰 $- I _ { \mathrm { h } }$ ，属六⽅晶系（图4.1），晶胞参数 $\dot { a } = b = 4 4 9 . 7 5 \mathrm { p m }$ $c = 7 3 2 . 2 4 \mathrm { p m } _ { \mathrm { c } }$ o4-1-1 设冰- $- I _ { \mathrm { h } }$ 中键⻆均为109.5°，计算氢键键⻓。

\subsection*{4-1}
-2 冰- $- I _ { \mathrm { h } }$ 理想模型中与每个氧原⼦相连的4个氧原⼦构成正四⾯体。计算实际晶体与理想模型的晶胞参数�的相对偏差（⽤百分数表⽰）。

\subsection*{4-2}
冰-XI 是冰 $- I _ { \mathrm { h } }$ 在低温下的平衡结构，可在低于72 K，常压下从KOH稀溶液中析出，其⽹状结构与冰 $- I _ { \mathrm { h } }$ 相似，但氢原⼦位置是有序的。冰-XI属正交晶系，晶胞参数 $\mathrm { \Delta } a = 4 5 0 . 1 9 \mathrm { p m }$ $b = 7 7 9 . 7 8 \mathrm { p m }$ $c = 7 3 2 . 8 0$ pm，计算该晶胞中⽔分⼦的个数。

\subsection*{4-3}
冰- $\boldsymbol { \cdot } \boldsymbol { I } _ { \mathrm { c } } .$ 是⼀种亚稳态，在− $^ { - 8 0 ~ ^ { \circ } C }$ 由⽔蒸汽凝结⽽成，属⽴⽅晶系，晶体内所有氢键键⻓均相等，晶胞参数�$= 6 3 5 . 8 0 \mathrm { p m } _ { \mathrm { c } }$ 。计算氢键键⻓并给出该晶体的点阵类型（不考虑氢原⼦，须写出简要的分析过程)。

（图：）

（图：）
五⻆⼗⼆⾯体

（图：）
⼗四⾯体
图 4.1 ⼩球为O原⼦
图 4.2

\subsection*{4-4}
在⾼压下，可得到⽴⽅晶系的冰- $\cdot \mathrm { X } ,$ ，晶胞参数 $\dot { a } = 2 7 8 . 5 0 \mathrm { p m }$ 。在冰-X中，所有氢键键⻓均相等，且氢原⼦在相邻两个氧原⼦的中点。求冰-X的密度与冰- ${ } ^ { - } I _ { \mathrm { c } }$ 密度的⽐值并确定其点阵类型（须写出简要的分析过程）。

\subsection*{4-5}
⽔与⼀些⽓体形成笼状化合物晶体。⽓体A与⽔形成的晶体由五⻆⼗⼆⾯体和⼗四⾯体两种⽔笼（图4.2）构成，每个⽔笼中⼼均包含⼀个A分⼦（视为球形），X射线衍射表明该晶体属⽴⽅晶系，晶胞中⼼为对称中⼼，体对⻆线为3次旋转轴。已知其中两个A的分数坐标分别为0, 0, 0和0, 1/2, 1/4。写出该晶胞所有A的分数坐标及相应⽔笼的类型；不考虑氢原⼦，写出该晶体的点阵类型（须写出简要的分析过程）。
